Under explicit Markov assumptions, the expected runs of a batting order can be computed exactly and
cheaply, which turns batting-order optimization from a noisy search into an exhaustive one with a
certificate. With an empirical-Bayes plate-appearance model that uses only pregame information,
the model's run predictions are calibrated on two held-out seasons, and the simple model predicts
plate appearances about as well as a gradient-boosted classifier given the same inputs.

The substantive conclusion is modest and, we think, robust: in this model the batting orders that
MLB teams actually used in 2025--2026 were close to optimal (median rank above \EfourbPctMedian{}
of all orders), and reordering would add about \EfourbUpliftMean{} expected runs per game, roughly
\EfourbUpliftSeason{} runs over a season. On average, the best and worst possible orders of the
same nine hitters differ by less than a tenth of a run per game (\EfourbSpread{}). For the 2026 World Series Game 5 the
estimated gains were \RefUpliftAway{} (Milwaukee) and \RefUpliftHome{} (Tampa Bay) runs, with
parameter uncertainty of the same size, and the identity of the single best order is not
statistically determined. Both orders were near-optimal: for Tampa Bay the model's preferred order is better
with high posterior probability (\RefPBetterHome{} of draws), but by only \RefUpliftHome{} runs per
game, and for Milwaukee the gain's interval includes zero. The actual result of the game cannot speak
to either question.

What remains uncertain is everything outside the model: base-running speed, mid-inning pitching
changes, individual platoon effects and in-game substitutions. Each of these could change the
ranking of nearly tied orders; we expect none of them to change the order of magnitude of the
effect. The method itself is the paper's practical contribution: a reproducible, leakage-free
pipeline and an exact optimizer, with an interactive public interface that runs the same
computation in the browser.
